% ═════════════════════════════════════════════════════════════════════════════
\section{Other Approaches and Related Chapters}
\label{sec:other}
% ═════════════════════════════════════════════════════════════════════════════

The techniques of Sections~\ref{sec:rmc} to~\ref{sec:views} are part of a
larger body of work. In this section, we relate them to other approaches, and
we point to the chapters of this handbook in which related topics are treated
in depth. More detailed surveys of parameterized verification can be found
in~\cite{AST18,BJKKRVW15,Esp14,ZP04}.

\subsection{Cutoffs}

A \emph{cutoff} for a class of systems is a number~$k$ such that a property
holds for all instances of a system whenever it holds for all instances of size
at most~$k$. A cutoff theorem thus reduces parameterized verification to
finitely many finite-state problems. Emerson and Namjoshi~\cite{EN95} showed
that, in rings in which a single token without data is passed between the
processes, properties that refer to a fixed number of processes have a small
cutoff that does not depend on the process template. Emerson and
Kahlon~\cite{EK00} considered \emph{guarded protocols}, in which a transition is
guarded by a \emph{conjunctive} condition, requiring all other processes to be
in given states, or by a \emph{disjunctive} condition, requiring at least one
of them to be. In the notation of Section~\ref{sec:prelim}, these are
universal and existential conditions with range~$\neq$. For disjunctive
guards, each guard needs only one witness, so a small instance that keeps one
process in each local state that occurs in a run can simulate the run. For
conjunctive guards, removing processes can only make the guards easier to
satisfy, so the projection of a run onto a subset of the processes is a run of
a smaller instance. The same two arguments explain why view abstraction only
needs configurations of size $k+1$ to compute the abstract post-image
(Section~\ref{sec:views}). Cutoff results have been extended to other
topologies~\cite{CTTV04,AJKR14}, to fairness and deadlock
detection~\cite{AJK16}, and to cache coherence protocols~\cite{EK03}. Static
cutoffs hold only for specific classes of systems, and they may be large.
\emph{Dynamic} cutoff methods, such as view abstraction and the method of
Kaiser, Kroening, and Wahl~\cite{KKW10} for programs with an unbounded number
of threads, find the cutoff during the analysis.

\subsection{Abstraction and Compositional Reasoning}

Several approaches keep a small number of processes precise and summarize the
others, in the same way as view abstraction keeps $k$ processes.
\emph{$(0,1,\infty)$-counter abstraction}~\cite{PXZ02} records, for each local
state, whether zero, one, or many processes occupy it, and adds fairness
conditions so that liveness properties can be proved. \emph{Environment
abstraction}~\cite{CTV06,CTV08} describes the system from the point of view of
one reference process, by predicates that relate the other processes to it.
\emph{Indexed predicates}~\cite{LB07} lift predicate abstraction to quantified
invariants. For cache coherence protocols, compositional model
checking~\cite{McM01} and the CMP method~\cite{CMP04,TT08} keep a few
processes concrete and replace the others by an abstract process, whose
behavior is constrained by noninterference lemmas. \emph{Thread-modular}
reasoning verifies one thread against an assumption on the others~\cite{FQ03},
and its precision can be increased by reasoning about several threads at a
time~\cite{HMP17,FKP15}, much as view abstraction gains precision by
increasing~$k$. The network
invariants of Section~\ref{sec:early} are a compositional method of the same
kind: a single process abstracts any number of copies of the template, and it
can be found by induction~\cite{KM89,WL89}, computed
automatically~\cite{LHR97}, or established deductively~\cite{KPSZ02}.
Abstraction and abstraction refinement are the subject of
\chapref{13}, predicate abstraction of \chapref{15}, and compositional
reasoning of \chapref{12}.

\subsection{Deductive and SMT-Based Methods}

Deductive methods prove that a quantified formula is an inductive invariant
for an arbitrary number of processes. The approach goes back to the proof method of
Owicki and Gries~\cite{OG76}, and it is supported by languages and tools such
as TLA$^+$~\cite{Lam94,Lam02,CDLMRV12}, which has been applied to industrial
protocols~\cite{NRZBMD15}. The main difficulty is to find the invariant, and
much work aims at inferring it automatically. The \emph{invisible invariants}
method~\cite{PRZ01,APRXZ01} generalizes the reachable configurations of a
small instance to a quantified candidate invariant, in a similar way as view
abstraction generalizes the views of small instances. In \emph{array-based
systems}~\cite{GR10,GS10} and in the Cubicle model checker~\cite{CGKMZ12,CGKMZ13},
backward reachability is carried out on existentially quantified formulas with
an SMT solver~\cite{BSST09}, and universal guards are handled by an
approximation closely related to monotonic abstraction. The IC3/PDR
algorithm~\cite{Bra11,EMB11} has been extended to infer universally quantified
invariants~\cite{KBIRS17,GSM16}. The Ivy tool~\cite{PMPSS16} restricts
verification conditions to a decidable fragment of first-order logic, in which
consensus protocols such as Paxos have been verified~\cite{PLSS17}, and the
inference of invariants for distributed protocols has been
automated~\cite{MGJKKS19,KPIA20}. For fault-tolerant distributed algorithms,
\emph{threshold automata} reduce parameterized verification to bounded model
checking with an SMT solver~\cite{JKSVW13,KVW17,KLVW17}. Satisfiability modulo
theories is the subject of \chapref{11}, the combination of model checking and
automated deduction of \chapref{20}, and automated deductive program
verification of \chapref{35}.

\subsection{Related Models and Applications}

\paragraph{Well-structured transition systems.}
The counting abstraction of Section~\ref{sec:early} turns systems of anonymous
finite-state processes into Petri nets~\cite{GerSis92}, or into broadcast
protocols when processes can broadcast~\cite{EFM99,Del00}. These models, the
backward reachability analysis of Section~\ref{sec:wqo}, and its extensions are
the subject of \chapref{33}. Monotonic abstraction (Section~\ref{sec:monotonic})
can be seen as a way to apply this theory to systems that are not
well-structured, and view abstraction (Section~\ref{sec:views}) gives a forward
alternative to backward reachability for well-structured systems.

\paragraph{Weak memory models.}
Synchronization primitives, such as locks, must be correct for any number of
threads under the memory model of the architecture, in which the reordering of
memory operations is described by unbounded buffers~\cite{OSS09,SSONM10,AMT14}. For a fixed number of threads,
safety verification under Total Store Order (TSO) is decidable by a reduction
to lossy channel systems~\cite{ABBM10}. A \emph{dual} semantics, in which
threads read from load buffers instead of writing to store buffers, makes the
system monotonic with respect to a WQO, which gave the first decidability
result for parameterized verification under TSO~\cite{AABN16,AABN18}. The
problem is \textsc{PSpace}-complete when the threads do not use atomic
read-modify-write operations~\cite{AAR20}, and parameterized verification under
the release-acquire semantics of C11 is also \textsc{PSpace}-complete under
restrictions on the use of compare-and-swap operations~\cite{KGMC22}. For the
Intel x86 model with non-temporal writes~\cite{RMV22}, verification is
undecidable in general, but it becomes decidable when each thread alternates a
bounded number of times between TSO and PSO behavior~\cite{AABKS24}, also for
an arbitrary number of threads, through well-structured systems based on
\emph{better quasi-orderings}~\cite{AABKS26,Mar94,Abd10}. Weak memory models are
treated in \chapref{36}.

\paragraph{Parameterized quantum circuits.}
Parameterized verification has recently been applied to quantum programs, in
which the parameter is the number of qubits. Quantum algorithms such as
Grover's search~\cite{Gro96} are described by programs that generate a circuit
for each number of qubits, and their correctness relates input states to
output states for all sizes. One approach~\cite{AutoQPara} represents infinite
families of quantum states by \emph{synchronized weighted tree automata} and
quantum gates by transducers, in a way similar to regular model checking
(Section~\ref{sec:rmc}), and reduces verification to decidable inclusion and
equivalence problems for these automata. It builds on automata-based
verification of quantum circuits of fixed
size~\cite{CCLLTY23,CCLLT23,LSTA25,AutoQ20}. Verification of quantum systems,
including deductive approaches~\cite{Yin11,Qbricks,HRHWH21,Amy19}, is the
subject of \chapref{45}.

\paragraph{Other chapters.}
The automata-theoretic background of regular model checking is presented in
\chapref{4}, and model checking of concurrent programs is the subject of
\chapref{18}.
